\chapter{Yield Curve Representations: A Literature Review}
\label{ch:model-background}

\textcolor{red}{Modelling the term structure of interest rates has been of broad interest in fixed income research.}  The yield curve is not observed directly but must be inferred from market observables such as the swap rates introduced in Chapter~\ref{chap:datadescription}, and since it is observed across many maturities, a direct representation quickly becomes high-dimensional. The central question is therefore whether it can be transformed into a smaller-dimensional representation while retaining as much information as possible \citep{Skovmand2025tsm}.

Classical approaches address this through representations that are simple and tractable. The Nelson--Siegel family \citep{NelsonSiegel1987} is perhaps the most well-known example, imposing a fixed parametric form on the curve. An alternative is to make no parametric assumption and instead apply linear dimension-reduction techniques such as principal component analysis \citep{Skovmand2025tsm}. Both approaches share the common goal of representing the yield curve with as few factors as possible, though Nelson-Siegel fixes the functional form a priori while PCA determines the factors empirically from the data.

However, Nelson--Siegel is constrained by its fixed functional form, while PCA is limited to linear combinations of the observed rates. Neural network-based approaches address both of these limitations by learning flexible, nonlinear representations directly from data. lets\citet{Kondratyev2018} was among the first to demonstrate that neural networks outperform classical dimensionality reduction techniques such as PCA in describing yield curve dynamics. \citet{Sokol2022} extended this by proposing a variational autoencoder as a data-driven replacement for classical parametric factor models, requiring no a priori choice of functional form for the curve. He further shows that training on data from multiple currencies simultaneously improves extrapolation beyond what a single currency has historically experienced.\footnote{In this context, extrapolation refers to the model being asked to produce output for curve shapes that have not been observed for a given currency during training. As \citet{Sokol2022} notes, any mapping from the yield curve to model state variables must be capable of describing such curve shapes in order to be useful for model construction.} By encoding all curves into a shared latent space, curve shapes unseen for one currency are supported by observations from others.

Such data-driven approaches, however, do not explicitly enforce no-arbitrage conditions on the generated curves. \citet{Andreasen2023} builds on this line of work with this in mind, interpreting the decoder output as zero-coupon bond prices rather than swap rates directly, which allows the stochastic dynamics of the encoded factors to be inferred, yielding a model that is close to dynamically arbitrage-free.

\citet{MatinWanPoulsen2025} take this a step further by imposing the no-arbitrage condition as a hard constraint at the decoding stage, rather than the soft regularisation term used by \citet{Andreasen2023}, producing yield curves that are exactly arbitrage-free by construction. The present thesis builds on this model, which we describe in detail in 
the following chapter.







